\documentclass[10pt,a4paper]{amsart}
\usepackage[margin=19mm]{geometry}
\usepackage{amsmath,amssymb,mathtools}
\usepackage[scr=rsfs,cal=euler]{mathalfa}
\usepackage{xfrac}
\usepackage[unicode,hidelinks,bookmarks=false]{hyperref}
\newcommand{\ie}{i.e.\ }
\newcommand{\cf}{cf.\ }
\newcommand{\resp}{respectively\ }
\newcommand{\Subsection}{\subsection}
\providecommand{\nobreakdash}{\nobreak\hbox{-}}
\setlength{\emergencystretch}{2em}
\multlinegap0pt
\usepackage{amssymb}
\usepackage{xcolor}
\newtheorem{lemma}{Lemma}
\def\DD{{\mathcal{D}}}
\def\a{\alpha}
\title{Pointwise convergence of the non-linear Fourier transform}
\author{Alexei Poltoratski}
\date{Source: arXiv:2103.13349v14 (CC BY 4.0)}
\begin{document}
\maketitle

\noindent\emph{Source and licence.} Pointwise convergence of the non-linear Fourier transform. Version arXiv:2103.13349v14 (CC BY 4.0), distributed by arXiv under the Creative Commons Attribution 4.0 International licence, http://creativecommons.org/licenses/by/4.0/, as recorded in the arXiv licence metadata for this exact version and shown on the abstract page https://arxiv.org/abs/2103.13349v14. Full licence notice: \texttt{licenses/arxiv-2103.13349-CC-BY-4.0.txt}.\par\smallskip

\noindent\textit{Editorial scope.} Counted unit: the article's lemma lemAppendix (source0.tex lines 2570--2589). The article's complete original proof of it is delivered with it (source0.tex lines 2591--2720, one \texttt{proof} environment of 130 lines). No mathematics is written, repaired or re-derived: the statement and the proof are the article's own lines, in the article's own notation, and no retained line is substituted or re-flowed.  No further source-local context is needed: every cross-reference of every retained line resolves inside the two delivered spans.  Nothing was omitted from any retained span, so the package carries no omission record.  No retained line contains a citation, so the wrapper carries no bibliography and no bibliography entry.  No retained line contains a figure environment, an image-inclusion command or drawing code, so no image is delivered and the package declares no asset dependency.  Editorial formatting only, in addition: an AMS \texttt{amsart} preamble that restates the article's own declarations and macros used by the retained lines, namely the environments \texttt{lemma} on separate counters, together with the standard packages those lines need.  The retained environments print in this document as Lemma 1 in order of appearance, whereas the article prints the same items with the numbering of its own full document.  The complete native source of the article as distributed in the arXiv e-print for this exact version is bundled unchanged as \texttt{sources/109525/source0.tex}.
\begin{lemma}\label{lemAppendix}
	Let $v>0$ be fixed. Let $\Delta$ be  a small positive constant and let 
	$\Delta_1>0$ and $\Delta_2>0$ be such that
	\[
	\left|\frac{\sinh \Delta_1}{\sinh(v+\Delta_1)}
	-\frac{\sinh \Delta}{\sinh(v-\Delta)}\right| <\Delta^3
	\]and \[
	\left|\frac{\sinh \Delta_2}{\sinh(v+\Delta_1+\Delta_2)}
	-\frac{\sinh \Delta}{\sinh(v-2\Delta)}\right| <\Delta^3.
	\]
	Then there exist constants $c=c(v)>0$ and $\Delta_0=\Delta_0(v)>0$ such that
	for all $0<\Delta<\Delta_0$,
	\[
	\frac{\sinh(2\Delta)}{\sinh(v-2\Delta)}
	-
	\frac{\sinh(\Delta_1+\Delta_2)}{\sinh(v+\Delta_1+\Delta_2)}
	\;>\;
	c\,\Delta^2 .
	\]
\end{lemma}

\begin{proof}
	Put
	\[
	F_v(x)=\frac{\sinh x}{\sinh(v+x)}, 
	\qquad s_v=\sinh v,\quad c_v=\cosh v.
	\]
	
	\medskip
	\noindent\textbf{Expansion of $F_v(x)$ at $x=0$.}  
	We use
	\[
	\sinh(v+x)=s_v + c_v x + \tfrac12 s_v x^2 + O(x^3), 
	\qquad
	\sinh x = x + \tfrac{1}{6}x^3 + O(x^5).
	\]
	Thus
	\[
	F_v(x)
	= \frac{x}{s_v} - \frac{c_v}{s_v^2}x^2 + O(x^3).
	\]
	Consequently,
	\[
	\frac{\sinh \Delta}{\sinh(v-\Delta)}
	= \frac{\Delta}{s_v} + \frac{c_v}{s_v^2}\Delta^2 + O(\Delta^3),
	\qquad
	\frac{\sinh (2\Delta)}{\sinh(v-2\Delta)}
	= \frac{2\Delta}{s_v} + \frac{4c_v}{s_v^2}\Delta^2 + O(\Delta^3).
	\]
	
	\medskip
	\noindent\textbf{Expansion of $\Delta_1$.}
	To expand $\DD_1$ in terms of $\DD$, let us write
	\[
	\Delta_1=\Delta+\alpha_1 \Delta^2+O(\Delta^3).
	\]
	Since
	\[
	F_v(\Delta_1)=\frac{\sinh\Delta}{\sinh(v-\Delta)} + O(\Delta^3),
	\]
	matching coefficients of $\Delta^2$ gives
	\[
	\frac{\alpha_1}{s_v}-\frac{c_v}{s_v^2}= \frac{c_v}{s_v^2},
	\]
	so
	\[
	\alpha_1=\frac{2c_v}{s_v},
	\qquad
	\Delta_1=\Delta+\frac{2c_v}{s_v}\Delta^2+O(\Delta^3).
	\]
	
	\medskip
	\noindent\textbf{Expansion of $\Delta_2$.} 
	To obtain an expansion for $\DD_2$ we can write
	\[
	\Delta_2=\Delta+\alpha_2\Delta^2+O(\Delta^3),\qquad
	u:=\Delta_1+\Delta_2.
	\]
	Using the expansion of $\Delta_1$ we have
	\[
	u=2\Delta+(\alpha_1+\alpha_2)\Delta^2+O(\Delta^3).
	\]
	The (approximate) defining equation is
	\[
	\frac{\sinh\Delta_2}{\sinh(v+u)}
	= \frac{\sinh\Delta}{\sinh(v-2\Delta)} + O(\Delta^3).
	\]
	
	The numerator expands as
	\[
	\sinh\Delta_2=\Delta+\alpha_2\Delta^2+O(\Delta^3).
	\]
	The denominator satisfies
	\[
	\sinh(v+u)=s_v + 2c_v\Delta 
	+ \bigl(c_v(\alpha_1+\alpha_2)+2s_v\bigr)\Delta^2 + O(\Delta^3).
	\]
	Using the reciprocal expansion and keeping terms up to $\Delta^2$ gives
	\[
	\frac{\sinh\Delta_2}{\sinh(v+u)}
	= \frac{\Delta}{s_v}
	+ \Bigl(-\frac{2c_v}{s_v^2}+\frac{\alpha_2}{s_v}\Bigr)\Delta^2
	+ O(\Delta^3).
	\]
	The right-hand side expands as
	\[
	\frac{\sinh\Delta}{\sinh(v-2\Delta)}
	= \frac{\Delta}{s_v} + \frac{2c_v}{s_v^2}\Delta^2 + O(\Delta^3).
	\]
	Matching the coefficients of $\Delta^2$ gives
	\[
	-\frac{2c_v}{s_v^2}+\frac{\alpha_2}{s_v}
	= \frac{2c_v}{s_v^2},
	\qquad
	\text{so}\qquad
	\alpha_2=\frac{4c_v}{s_v}\qquad
	\text{and}\qquad
\a_1+	\alpha_2=\frac{6c_v}{s_v}.
	\]
	Thus
	\[
	\Delta_2=\Delta+\frac{4c_v}{s_v}\Delta^2+O(\Delta^3)\quad\text{and}\quad
	u=\Delta_1+\Delta_2=2\Delta+\frac{6c_v}{s_v}\Delta^2+O(\Delta^3).
	\]
	
	\medskip
	\noindent\textbf{Final expansion.}
	Using the expansion of $F_v$,
	\[
	\frac{\sinh(u)}{\sinh(v+u)}
	= \frac{2\Delta}{s_v}+\frac{2c_v}{s_v^2}\Delta^2 + O(\Delta^3).
	\]
	Subtracting from the expansion of 
	\(\dfrac{\sinh(2\Delta)}{\sinh(v-2\Delta)}\) gives
	\[
	\frac{\sinh(2\Delta)}{\sinh(v-2\Delta)}
	-
	\frac{\sinh(\Delta_1+\Delta_2)}{\sinh(v+\Delta_1+\Delta_2)}
	= \frac{2c_v}{s_v^2}\Delta^2 + O(\Delta^3),
	\]
	which implies the statement.
%	Since $2c_v/s_v^2>0$, for sufficiently small $\Delta$ the error term is 
%	bounded below by $-\tfrac12(2c_v/s_v^2)\Delta^2$, and therefore
%	\[
%	\frac{\sinh(2\Delta)}{\sinh(v-2\Delta)}
%	-
%	\frac{\sinh(\Delta_1+\Delta_2)}{\sinh(v+\Delta_1+\Delta_2)}
%	> \frac{c_v}{s_v^2}\,\Delta^2.
%	\]
%	The lemma follows with $c=c_v/s_v^2>0$.
\end{proof}

\end{document}
